% =============================================================================
% Chapter 7 — The Wavetable Synthesizer ($A140–$A1DF)
% =============================================================================
\chapter{The Wavetable Synthesizer}

\epigraph{``Any sufficiently advanced synthesizer is indistinguishable from an
orchestra.''}{\textit{--- With apologies to Arthur C.\ Clarke}}

\section*{Introduction}

The Wavetable Synthesizer --- WTS --- is an 8-voice sample-based sound engine
that complements the SID chips' analog-style synthesis with realistic instrument
timbres.  It loads standard SF2 soundfont files from the host filesystem,
keeping sample data in host memory, and renders stereo PCM16 audio at
44\,100\,Hz through OpenAL.

Each voice provides independent control over note pitch, velocity, instrument
selection, volume, stereo panning, and pitch bend.  Two global effects ---
reverb (Freeverb algorithm) and chorus (modulated delay) --- process the mixed
output before the master volume stage.

The WTS occupies the address range \$A140--\$A1DF, divided into three regions:
8 voice register blocks (\$A140--\$A17F), global control registers
(\$A180--\$A187), and the instrument enumeration interface (\$A1A0--\$A1DF).


% ═══════════════════════════════════════════════════════════════════════════════
\section{Voice Registers (\$A140--\$A17F)}
\label{sec:wts-voices}
% ═══════════════════════════════════════════════════════════════════════════════

Eight voices, eight bytes each.  The base address for voice~$N$ is
\$A140~+~$N \times 8$.

\begin{center}
\small
\begin{tabular}{@{}clcll@{}}
\toprule
\textbf{Voice} & \textbf{Base} & \textbf{End} & \textbf{Range} \\
\midrule
0 & \texttt{\$A140} & \texttt{\$A147} & Voice 0 \\
1 & \texttt{\$A148} & \texttt{\$A14F} & Voice 1 \\
2 & \texttt{\$A150} & \texttt{\$A157} & Voice 2 \\
3 & \texttt{\$A158} & \texttt{\$A15F} & Voice 3 \\
4 & \texttt{\$A160} & \texttt{\$A167} & Voice 4 \\
5 & \texttt{\$A168} & \texttt{\$A16F} & Voice 5 \\
6 & \texttt{\$A170} & \texttt{\$A177} & Voice 6 \\
7 & \texttt{\$A178} & \texttt{\$A17F} & Voice 7 \\
\bottomrule
\end{tabular}
\end{center}

The register offsets within each voice block are identical.  In the entries
below, the address shown is for voice~0; add $N \times 8$ for voice~$N$.

% --- +0 WtsVoiceNote ----------------------------------------------------------

\mementry{A140}{WtsVoiceNote}{W}{---}

MIDI note number.  Writing a value from 1 to 127 triggers a note-on at that
pitch (60~=~middle~C).  Writing~0 triggers a note-off, releasing the voice into
its envelope's release phase.

Set the velocity and instrument \emph{before} writing the note number, as the
write to this register is what triggers playback.

\begin{lstlisting}
POKE $A140, 60 : REM note on: middle C
POKE $A140, 0  : REM note off
\end{lstlisting}

% --- +1 WtsVoiceVelocity ------------------------------------------------------

\mementry{A141}{WtsVoiceVelocity}{W}{---}

Note velocity (0--127).  Affects the volume and, depending on the soundfont,
the timbre of the note.  Set this before writing the note number.

\begin{lstlisting}
POKE $A141, 100 : REM velocity 100 (moderately loud)
\end{lstlisting}

% --- +2 WtsVoiceInstrument ----------------------------------------------------

\mementry{A142}{WtsVoiceInstrument}{W}{---}

Instrument index (0--255).  Selects from the currently loaded soundfont.  The
number of available instruments can be read from \addr{A184}.  Set this before
playing notes.

\begin{lstlisting}
POKE $A142, 0   : REM instrument 0 (typically piano)
POKE $A142, 40  : REM instrument 40 (typically violin)
\end{lstlisting}

% --- +3 WtsVoiceVolume --------------------------------------------------------

\mementry{A143}{WtsVoiceVolume}{R/W}{255}

Per-voice volume (0--255).  Independent of velocity --- this acts as a
channel-level volume control.  The final output level is the product of
velocity, voice volume, and master volume.

\begin{lstlisting}
POKE $A143, 128 : REM half volume on voice 0
\end{lstlisting}

% --- +4 WtsVoicePanning -------------------------------------------------------

\mementry{A144}{WtsVoicePanning}{R/W}{128}

Stereo position.  0~=~hard left, 128~=~centre, 255~=~hard right.

\begin{lstlisting}
POKE $A144, 0   : REM voice 0 hard left
POKE $A14C, 255 : REM voice 1 hard right
\end{lstlisting}

% --- +5 WtsVoicePitchBendLo ---------------------------------------------------

\mementry{A145}{WtsVoicePitchBendLo}{W}{0}

Pitch bend low byte.  Combined with the high byte (\addr{A146}) to form a
16-bit pitch bend value.

% --- +6 WtsVoicePitchBendHi ---------------------------------------------------

\mementry{A146}{WtsVoicePitchBendHi}{W}{\$80}

Pitch bend high byte.  The centre (no bend) value is \$80.  Values above \$80
bend the pitch upward; values below bend downward.  The bend range depends on
the soundfont preset but is typically $\pm$2 semitones.

\begin{lstlisting}
10 REM bend voice 0 up one step
20 POKE $A145, 0    : REM low byte
30 POKE $A146, $C0  : REM high byte (above centre)
\end{lstlisting}

% --- +7 WtsVoiceStatus --------------------------------------------------------

\mementry{A147}{WtsVoiceStatus}{R}{0}

Voice status flags.

\bitfield{---}{---}{---}{---}{---}{---}{Rel}{Act}

\begin{description}
  \item[Bit 0 --- Active]\hfill\\
    Set while the voice is producing sound (note on or sustaining).
  \item[Bit 1 --- Releasing]\hfill\\
    Set when the voice is in the release phase of its envelope (after note off).
    Cleared when the release completes and the voice becomes silent.
\end{description}

\begin{lstlisting}
10 REM wait for voice 0 to finish releasing
20 WAIT $A147, 3, 3  : REM wait until bits 0+1 both clear
\end{lstlisting}


% ═══════════════════════════════════════════════════════════════════════════════
\section{Global Registers (\$A180--\$A187)}
\label{sec:wts-global}
% ═══════════════════════════════════════════════════════════════════════════════

% --- $A180 WtsReverbLevel -----------------------------------------------------

\mementry{A180}{WtsReverbLevel}{R/W}{80}

Reverb wet level.  0~=~completely dry (no reverb), 255~=~maximum reverb.  The
reverb uses a Freeverb algorithm with multiple comb and allpass filters.

\begin{lstlisting}
POKE $A180, 120 : REM generous reverb
POKE $A180, 0   : REM reverb off
\end{lstlisting}

% --- $A181 WtsChorusLevel -----------------------------------------------------

\mementry{A181}{WtsChorusLevel}{R/W}{40}

Chorus depth.  0~=~no chorus effect, 255~=~maximum depth.  The chorus uses a
modulated delay line to create a thickening, ensemble-like effect.

\begin{lstlisting}
POKE $A181, 80 : REM moderate chorus
\end{lstlisting}

% --- $A182 WtsMasterVolume ----------------------------------------------------

\mementry{A182}{WtsMasterVolume}{R/W}{255}

Master output volume for all WTS voices.  Scales the final mixed output before
it is sent to the audio device.  0~=~silent, 255~=~full volume.

\begin{lstlisting}
POKE $A182, 200 : REM slightly reduced master volume
\end{lstlisting}

% --- $A183 WtsSoundfontStatus -------------------------------------------------

\mementry{A183}{WtsSoundfontStatus}{R}{0}

Soundfont loading status:

\begin{center}
\small
\begin{tabular}{@{}cl@{}}
\toprule
\textbf{Value} & \textbf{Meaning} \\
\midrule
\$00 & No soundfont loaded \\
\$01 & Ready --- soundfont loaded and playable \\
\$02 & Loading --- soundfont is being loaded (wait before playing) \\
\$FF & Error --- soundfont failed to load \\
\bottomrule
\end{tabular}
\end{center}

\begin{lstlisting}
10 REM wait for soundfont to finish loading
20 WAIT $A183, 1   : REM wait for bit 0 (ready)
30 PRINT "SOUNDFONT READY"
\end{lstlisting}

% --- $A184 WtsInstrumentCount -------------------------------------------------

\mementry{A184}{WtsInstrumentCount}{R}{0}

Number of instruments available in the currently loaded soundfont.  Returns 0
if no soundfont is loaded.

\begin{lstlisting}
PRINT "INSTRUMENTS:"; PEEK($A184)
\end{lstlisting}

% --- $A185 WtsCommand ---------------------------------------------------------

\mementry{A185}{WtsCommand}{W}{---}

Command register.  Writing a command byte executes it immediately.

\cmdentry{01}{AllNotesOff}{Silence all voices}

Immediately stops all playing notes on all eight voices.

\cmdentry{02}{ResetEffects}{Reset reverb and chorus}

Resets reverb and chorus levels to their defaults (80 and 40 respectively) and
clears any residual audio in the effect buffers.

% --- $A186 WtsCommandArg -----------------------------------------------------

\mementry{A186}{WtsCommandArg}{W}{---}

Command argument byte.  Currently unused; reserved for future commands.

% --- $A187 WtsActiveVoices ----------------------------------------------------

\mementry{A187}{WtsActiveVoices}{R}{0}

Bitmask of currently active voices.  Bit~$N$ is set if voice~$N$ is producing
sound (either playing or releasing).

\bitfield{V7}{V6}{V5}{V4}{V3}{V2}{V1}{V0}

\begin{lstlisting}
10 REM check which voices are active
20 A = PEEK($A187)
30 FOR I=0 TO 7
40   IF A AND (2^I) THEN PRINT "VOICE"; I; "ACTIVE"
50 NEXT I
\end{lstlisting}


% ═══════════════════════════════════════════════════════════════════════════════
\section{Instrument Enumeration (\$A1A0--\$A1DF)}
\label{sec:wts-enum}
% ═══════════════════════════════════════════════════════════════════════════════

The enumeration interface allows programs to discover which instruments are
available in the loaded soundfont.  Write an instrument index to \addr{A1A0}
and the name, bank, and program number are populated for reading.

% --- $A1A0 WtsEnumIndex -------------------------------------------------------

\mementry{A1A0}{WtsEnumIndex}{W}{---}

Write an instrument index (0 to \texttt{PEEK(\$A184)-1}) to populate the
enumeration buffer.

% --- $A1A1 WtsEnumBank --------------------------------------------------------

\mementry{A1A1}{WtsEnumBank}{R}{---}

MIDI bank number of the selected instrument.

% --- $A1A2 WtsEnumProgram -----------------------------------------------------

\mementry{A1A2}{WtsEnumProgram}{R}{---}

MIDI program number of the selected instrument.

% --- $A1A3-$A1DF WtsEnumName --------------------------------------------------

\memrange{A1A3}{A1DF}{WtsEnumName}{R}{---}

Null-terminated instrument name string, up to 28 characters.

\begin{lstlisting}
10 REM list all instruments
20 N = PEEK($A184)
30 IF N=0 THEN PRINT "NO SOUNDFONT" : END
40 FOR I=0 TO N-1
50   POKE $A1A0, I          : REM select instrument
60   B = PEEK($A1A1)        : REM bank
70   P = PEEK($A1A2)        : REM program
80   N$ = ""
90   FOR J=0 TO 27
100    C = PEEK($A1A3+J)
110    IF C=0 THEN J=27 : GOTO 130
120    N$ = N$ + CHR$(C)
130  NEXT J
140  PRINT I; ": ["; B; ":"; P; "] "; N$
150 NEXT I
\end{lstlisting}


% ═══════════════════════════════════════════════════════════════════════════════
\section{Complete Examples}
\label{sec:wts-examples}
% ═══════════════════════════════════════════════════════════════════════════════

\subsection{Loading a Soundfont}

Soundfonts are loaded via the FileIoController at \texttt{\$B9A0}.
Command \$15 loads an SF2 file by name.

\begin{lstlisting}
10 REM load a soundfont via FileIoController
20 F$ = "GeneralUser.sf2" + CHR$(0)
30 FOR I=0 TO LEN(F$)-1
40   POKE $B9B0+I, ASC(MID$(F$,I+1,1))
50 NEXT I
60 POKE $B9A0, $15           : REM load soundfont command
70 WAIT $A183, 1             : REM wait for ready
80 PRINT PEEK($A184); " INSTRUMENTS LOADED"
\end{lstlisting}

\subsection{Playing a Note}

\begin{lstlisting}
10 REM play middle C on voice 0
20 BASE = $A140
30 POKE BASE+2, 0            : REM instrument 0 (piano)
40 POKE BASE+1, 100          : REM velocity 100
50 POKE BASE+0, 60           : REM note on: middle C
60 FOR I=1 TO 5000 : NEXT I  : REM wait
70 POKE BASE+0, 0            : REM note off
\end{lstlisting}

\subsection{Three-Note Chord with Panning}

\begin{lstlisting}
100 REM play a C major chord across the stereo field
110 FOR V=0 TO 2
120   B = $A140 + V*8
130   POKE B+2, 0             : REM piano
140   POKE B+3, 200           : REM volume
150   POKE B+1, 90            : REM velocity
160 NEXT V
170 POKE $A144, 64            : REM voice 0: left of centre
180 POKE $A14C, 128           : REM voice 1: centre
190 POKE $A154, 192           : REM voice 2: right of centre
200 REM play C-E-G
210 POKE $A140, 60            : REM C
220 POKE $A148, 64            : REM E
230 POKE $A150, 67            : REM G
240 FOR I=1 TO 10000 : NEXT I
250 REM release all
260 POKE $A140, 0
270 POKE $A148, 0
280 POKE $A150, 0
\end{lstlisting}

\subsection{Setting Up Reverb and Chorus}

\begin{lstlisting}
10 REM cathedral-style reverb with gentle chorus
20 POKE $A180, 180           : REM heavy reverb
30 POKE $A181, 60            : REM moderate chorus
40 POKE $A182, 220           : REM master volume
50 REM now any notes played will have these effects
\end{lstlisting}

\begin{tipbox}
The MML \texttt{MUSIC} command supports WTS voices.  Use \texttt{@I\textless n\textgreater}
to select a WTS instrument by GM program number, or \texttt{@D\textless n\textgreater}
to select a drum instrument.  For example, \texttt{MUSIC 6, "@I0 L4 CDEFG"} plays
a C major scale on WTS voice~0 using instrument~0.  This is often easier than
direct register manipulation for musical sequences.
\end{tipbox}
